\documentclass[12pt,a4paper]{article}

% ---- Packages ----
\usepackage[utf8]{inputenc}
\usepackage[T1]{fontenc}


\usepackage{amsmath,amssymb,amsthm}
\usepackage{mathtools}
\usepackage{geometry}
\usepackage{hyperref}
\usepackage{booktabs}
\usepackage{array}
\usepackage{graphicx}
\usepackage{xcolor}
\usepackage{mdframed}
\usepackage{enumitem}
\usepackage[numbers]{natbib}
\usepackage{doi}
\usepackage{setspace}

\geometry{margin=2.5cm}
\setstretch{1.12}

\hypersetup{
  colorlinks=true,
  linkcolor=blue!60!black,
  citecolor=blue!60!black,
  urlcolor=blue!60!black
}

% ---- Theorem environments ----
\theoremstyle{plain}
\newtheorem{theorem}{Theorem}[section]
\newtheorem{proposition}[theorem]{Proposition}
\newtheorem{lemma}[theorem]{Lemma}
\newtheorem{corollary}[theorem]{Corollary}
\theoremstyle{definition}
\newtheorem{definition}[theorem]{Definition}
\theoremstyle{remark}
\newtheorem{remark}[theorem]{Remark}

% ---- Erratum box ----
\newmdenv[
  linecolor=red!70!black,
  linewidth=1.4pt,
  backgroundcolor=red!5,
  frametitle={Mechanism Correction},
  frametitlefont=\bfseries\color{red!70!black},
  innertopmargin=8pt,
  innerbottommargin=8pt,
  skipabove=12pt,
  skipbelow=12pt
]{erratumbox}

% ---- Correction highlight box ----
\newmdenv[
  linecolor=blue!60!black,
  linewidth=1.2pt,
  backgroundcolor=blue!4,
  frametitle={Corrected Mechanism},
  frametitlefont=\bfseries\color{blue!70!black},
  innertopmargin=8pt,
  innerbottommargin=8pt,
  skipabove=10pt,
  skipbelow=10pt
]{correctionbox}

% ---- Macros ----
\newcommand{\Forb}{\mathcal{F}}
\newcommand{\CI}{\mathcal{C}_{\mathrm{I}}}
\newcommand{\CII}{\mathcal{C}_{\mathrm{II}}}
\newcommand{\mgap}[2]{g_{\min}(#1 \to #2)}

% ============================================================
\begin{document}
% ============================================================

\title{%
  \large\textsc{The Lemke Oliver--Soundararajan Bias as Surface Evidence}\\[4pt]
  \large\textsc{of Ternary Constitutional Structure}\\[12pt]
}

\author{%
  Khayyam Wakil\\[4pt]
  \small ARC Institute of Knowware\\
  \small Calgary, Alberta, Canada\\
  \small \href{mailto:the@knowware.institute}{\texttt{the@knowware.institute}}
}
\date{March 2026}

\maketitle

% ---- Prefatory note ----
\noindent\emph{This paper was not planned. On the Friday afternoon before
submitting the Constitutional Forcing programme to peer review, the author read a
popular account of the Lemke Oliver--Soundararajan prime digit bias---a result that has
been described as ``deeply unsettling'' since its publication in 2016. The explanation
appeared immediately, in a single conversation, without calculation. Constitutional
Forcing was already in the room. The bias recognized the framework. This paper is the
formalization of what took approximately one afternoon to see.}

% ============================================================
\begin{abstract}
In 2016, Lemke Oliver and Soundararajan published a striking computational discovery:
the last digits of consecutive primes exhibit a strong repulsion bias that contradicts
na\"ive randomness assumptions and persists---though diminishing---across all number
bases tested. The authors correctly traced this to the Hardy--Littlewood prime
constellations conjecture, but could not identify the underlying mechanistic cause:
why the structure exists, why it holds across all bases, and why the diminishment rate
is precisely what it is. We demonstrate that the bias is not mysterious: it is the direct
base-10 projection of ternary constitutional pressure governing prime distribution.
Specifically, the observed digit-repulsion is a shadow of the mod~3 constraint under
which all primes (except 2 and~3) reside in remainder classes $\{1, 2\}$, and under
which all twin prime pairs satisfy $p \equiv 2 \pmod{3}$. The \emph{snail's pace} of
diminishment is identified as the Siegel--Walfi{}sz convergence rate of ternary class
equilibration---not an empirical curiosity, but a fundamental quantity in analytic number
theory appearing in a context where it was not previously recognized. This explanation
is mechanistic rather than descriptive: it does not merely reproduce the bias from the
Hardy--Littlewood conjecture, but identifies the structural reason the conjecture itself
encodes such behaviour.
\end{abstract}

\noindent\textbf{Keywords:} prime digit bias, ternary constitutional structure,
Lemke Oliver--Soundararajan, mod~3 constraints, minimum gap structure, twin primes,
parity barrier, Siegel--Walfi{}sz theorem, cascade moduli

\noindent\textbf{MSC2020:} 11N05 (primary), 11N13, 11A41, 94A17 (secondary)

% ============================================================
\section{Introduction: The Mystery That Wasn't}
\label{sec:intro}
% ============================================================

In 2016, the mathematical world was surprised by a result so elementary that any
undergraduate could verify it in an afternoon---and yet no one had checked.

If primes were distributed randomly among their possible terminal digits, each of the
four digits $\{1, 3, 7, 9\}$ should follow any other with probability approximately 25\%.
This is what randomness means. Lemke Oliver and Soundararajan tested the first billion
primes~\cite{LOS16}. The answer violated the assumption.

A prime ending in~9 is followed by a prime ending in~1 roughly 65\% more often than by
another prime ending in~9. Every prime actively avoids repeating its own terminal digit.
The bias holds across every number base tested. It diminishes---but at what the authors
themselves called a snail's pace.

They correctly traced its theoretical origin to the Hardy--Littlewood prime constellations
conjecture~\cite{HL23}. But this explanation is descriptive rather than mechanistic. It
says that the Hardy--Littlewood conjecture, if true, implies the bias. It does not explain:
\begin{enumerate}[label=(\roman*)]
  \item Why the structure exists in the first place;
  \item Why it holds across all number bases, not just base~10;
  \item Why the diminishment rate is precisely what it is;
  \item What the bias is a symptom of at a deeper structural level.
\end{enumerate}
We answer all four questions. The answer is the same in each case: ternary constitutional
structure.

\medskip
\noindent\textbf{Main result (informal).} \emph{The Lemke Oliver--Soundararajan bias is
the base-10 projection of mod~3 constitutional pressure on prime distribution. The
diagonal suppression (primes avoiding their own terminal digit) reflects the structure of
minimum possible prime gaps, which carry fixed ternary class signatures. The snail's
pace of diminishment is the Siegel--Walfi{}sz convergence rate of ternary class
equilibration. Both phenomena are transparent once viewed in the correct coordinate
system.}

\medskip
\noindent\textbf{Organisation.} Section~\ref{sec:findings} recapitulates the Lemke
Oliver--Soundararajan findings. Section~\ref{sec:framework} develops the ternary
constitutional framework. Section~\ref{sec:mechanism} derives the bias mechanistically
(with erratum preceding it). Section~\ref{sec:snail} identifies the convergence rate.
Section~\ref{sec:validation} discusses base-independence, the parity barrier, and relation
to the Constitutional Forcing programme. Section~\ref{sec:conclusion} concludes.

% ============================================================
\section{The Lemke Oliver--Soundararajan Findings}
\label{sec:findings}
% ============================================================

\subsection{The Observation}

Working with the first billion primes, Lemke Oliver and Soundararajan~\cite{LOS16}
constructed a transition matrix. For consecutive primes $p$ and~$q$ (the next prime
after~$p$), they measured the frequency with which the terminal digit of~$p$ is followed
by each possible terminal digit of~$q$. The terminal digits of primes greater than~5
are drawn from $\{1, 3, 7, 9\}$.

The uniform (random) prediction is that each transition occurs with probability
$1/4 = 25\%$. The observed transition matrix (approximate, first billion primes, base~10) is:

\begin{center}
\renewcommand{\arraystretch}{1.25}
\begin{tabular}{c|cccc}
\toprule
$p \to q$ & ends in 1 & ends in 3 & ends in 7 & ends in 9 \\
\midrule
ends in 1 & ${\sim}18\%$ & ${\sim}30\%$ & ${\sim}30\%$ & ${\sim}22\%$ \\
ends in 3 & ${\sim}38\%$ & ${\sim}18\%$ & ${\sim}24\%$ & ${\sim}20\%$ \\
ends in 7 & ${\sim}31\%$ & ${\sim}24\%$ & ${\sim}19\%$ & ${\sim}26\%$ \\
ends in 9 & ${\sim}41\%$ & ${\sim}20\%$ & ${\sim}22\%$ & ${\sim}17\%$ \\
\bottomrule
\end{tabular}
\end{center}

The headline finding: \emph{every prime actively avoids repeating its own terminal digit}.
The diagonal entries are all suppressed below 25\%. The effect is most pronounced for
primes ending in~1 (suppressed to~18\%) and~9 (suppressed to~17\%) and mildest for
primes ending in~7 (suppressed to~19\%).

\begin{figure}[htbp]
\centering
\includegraphics[width=\textwidth]{fig1_transition_matrices}
\caption{Prime terminal digit transition matrices. \emph{Left:} observed transition
frequencies from the first billion primes (Lemke Oliver--Soundararajan~\cite{LOS16}).
\emph{Centre:} min-gap ternary structural prediction (corrected mechanism; see
Section~\ref{sec:mechanism}). \emph{Right:} deviation of observed frequencies from the
uniform 25\% baseline. Gold boxes highlight the suppressed diagonal (same-digit
self-transitions). The hot/cold structure of the prediction matches the observed pattern:
transitions whose minimum possible gap is $\equiv 0 \pmod{3}$ are redistributed;
transitions with minimum gap~$2$ (gap class~II) are most enhanced.}
\label{fig:matrices}
\end{figure}

\subsection{Prior Explanation and Its Limits}

Lemke Oliver and Soundararajan traced the bias to the Hardy--Littlewood prime
constellations conjecture~\cite{HL23}---specifically, to the singular series that weights
the expected density of prime pairs with specified gap sizes. The singular series correctly
predicts the transition probabilities to good accuracy.

However, this explanation is descriptive. It says: if the Hardy--Littlewood conjecture
holds, then the bias follows. It does not identify why the conjecture itself encodes the
pattern, why the pattern is base-independent, or why the snail's pace is structurally
determined. The authors described the phenomenon as ``deeply unsettling'' and noted that
primes ``somehow `know' about future primes.'' This language---primes possessing
knowledge, the result being unsettling---signals a missing mechanistic explanation. We
supply it.

% ============================================================
\section{The Ternary Constitutional Framework}
\label{sec:framework}
% ============================================================

\subsection{Core Definitions}

We work with remainder classes modulo~3. Every integer $n$ belongs to exactly one of
three classes:
\begin{itemize}
  \item $n \equiv 0 \pmod{3}$: the \textbf{Forbidden Class} (divisible by~3);
  \item $n \equiv 1 \pmod{3}$: \textbf{Constitutional Class~I};
  \item $n \equiv 2 \pmod{3}$: \textbf{Constitutional Class~II}.
\end{itemize}

\begin{definition}[Constitutional Constraint]
\label{def:constraint}
A prime $p > 3$ is \emph{constitutionally constrained}: it must satisfy
$p \equiv 1$ or $p \equiv 2 \pmod{3}$. The Forbidden Class is closed to primes
greater than~3.
\end{definition}

This is elementary: any integer $n > 3$ with $n \equiv 0 \pmod{3}$ is divisible by~3
and hence composite. Therefore all primes $p > 3$ lie in
$\{p : p \equiv 1 \text{ or } 2 \pmod{3}\}$.

\subsection{The Twin Prime Constitutional Signature}

\begin{theorem}[Constitutional Signature, {\cite[Prop.~4.1]{Wakil26b}}]
\label{thm:signature}
Let $(p, p+2)$ be a twin prime pair with $p > 3$. Then
\[
  p \equiv 2 \pmod{3} \quad \text{and} \quad p + 2 \equiv 1 \pmod{3}.
\]
\end{theorem}

\begin{proof}
Among any three consecutive integers, exactly one is divisible by~3.
\begin{itemize}
  \item If $p \equiv 0 \pmod{3}$: then $p = 3$ (the only such prime), excluded.
  \item If $p \equiv 1 \pmod{3}$: then $p + 2 \equiv 0 \pmod{3}$, so $p + 2 > 3$ is
    composite. Contradiction.
  \item If $p \equiv 2 \pmod{3}$: then $p + 2 \equiv 1 \pmod{3}$. Neither $p$ nor
    $p + 2$ is divisible by~3. Both may be prime. \qedhere
\end{itemize}
\end{proof}

\begin{remark}
Theorem~\ref{thm:signature} is verified computationally for all 3{,}424{,}009 twin
prime pairs with $p \leq 10^9$, with zero exceptions (beyond the special pair $(3, 5)$).
\end{remark}

\subsection{The Ternary Equilibrium}

\begin{theorem}[Ternary Equilibrium]
\label{thm:equilibrium}
Let $\pi_1(N)$ and $\pi_2(N)$ denote the number of primes $p \leq N$ with
$p \equiv 1 \pmod{3}$ and $p \equiv 2 \pmod{3}$ respectively. Then
\[
  \frac{\pi_1(N)}{\pi(N)} \to \frac{1}{2} \quad \text{and} \quad
  \frac{\pi_2(N)}{\pi(N)} \to \frac{1}{2} \quad \text{as } N \to \infty.
\]
\end{theorem}

\begin{proof}
Dirichlet's theorem on primes in arithmetic progressions~\cite{Dir37}: since
$\gcd(1,3) = \gcd(2,3) = 1$, both residue classes modulo~3 contain infinitely many
primes with equal asymptotic density.
\end{proof}

% ============================================================
\section{Deriving the Lemke Oliver--Soundararajan Bias}
\label{sec:mechanism}
% ============================================================

\subsection{The Mechanism: Minimum Gap Ternary Structure}

\begin{correctionbox}
\noindent Each terminal-digit transition $(d_1 \to d_2)$ has a fixed
\textbf{minimum possible prime gap}, determined entirely by decimal arithmetic: the
smallest positive integer $g$ such that if $p$ ends in~$d_1$, then $p + g$ ends in~$d_2$.
The ternary class of this minimum gap governs whether the transition is enhanced,
suppressed, or redistributed.
\end{correctionbox}

\medskip

The minimum gap $\mgap{d_1}{d_2}$ is the smallest positive even integer $g$ such that
$(d_1 + g) \equiv d_2 \pmod{10}$. Since all prime gaps (for $p > 5$) are even, this is
determined by the decimal digit arithmetic of $d_1$ and $d_2$. The complete assignment
is given in Table~\ref{tab:mingaps}.

\begin{table}[htbp]
\centering
\caption{Minimum possible prime gap for each terminal-digit transition, with its ternary
class. The ternary class of the gap---not the ternary class of the terminal digit---is the
structural driver.}
\label{tab:mingaps}
\renewcommand{\arraystretch}{1.3}
\begin{tabular}{cccccc}
\toprule
Transition & Min gap $g$ & $g \pmod{3}$ & Ternary class & Category & Observed effect \\
\midrule
$9 \to 1$ & 2 & 2 & Class~II & gap-2 & Enhanced (${\sim}41\%$, $1.64\times$) \\
$7 \to 9$ & 2 & 2 & Class~II & gap-2 & Enhanced (${\sim}26\%$, $1.04\times$) \\
$1 \to 3$ & 2 & 2 & Class~II & gap-2 & Enhanced (${\sim}30\%$, $1.20\times$) \\
\midrule
$3 \to 7$ & 4 & 1 & Class~I  & gap-4 & Moderate (${\sim}24\%$, $0.96\times$) \\
$7 \to 1$ & 4 & 1 & Class~I  & gap-4 & Enhanced (${\sim}31\%$, $1.24\times$) \\
$9 \to 3$ & 4 & 1 & Class~I  & gap-4 & Moderate (${\sim}20\%$, $0.80\times$) \\
\midrule
$1 \to 7$ & 6 & 0 & Forbidden & gap-6 (redist.) & Elevated (${\sim}30\%$, $1.20\times$) \\
$3 \to 9$ & 6 & 0 & Forbidden & gap-6 (redist.) & Elevated (${\sim}20\%$, $0.80\times$) \\
$7 \to 3$ & 6 & 0 & Forbidden & gap-6 (redist.) & Moderate (${\sim}24\%$, $0.96\times$) \\
\midrule
$3 \to 1$ & 8 & 2 & Class~II & gap-8 & Enhanced (${\sim}38\%$, $1.52\times$) \\
$9 \to 7$ & 8 & 2 & Class~II & gap-8 & Moderate (${\sim}22\%$, $0.88\times$) \\
$1 \to 9$ & 8 & 2 & Class~II & gap-8 & Moderate (${\sim}22\%$, $0.88\times$) \\
\midrule
$1 \to 1$ & 10 & 1 & Class~I  & same-digit & Suppressed (${\sim}18\%$, $0.72\times$) \\
$3 \to 3$ & 10 & 1 & Class~I  & same-digit & Suppressed (${\sim}18\%$, $0.72\times$) \\
$7 \to 7$ & 10 & 1 & Class~I  & same-digit & Suppressed (${\sim}19\%$, $0.76\times$) \\
$9 \to 9$ & 10 & 1 & Class~I  & same-digit & Suppressed (${\sim}17\%$, $0.68\times$) \\
\bottomrule
\end{tabular}
\end{table}

\subsection{Three Structural Categories}

The 16 transitions partition naturally into three categories according to the ternary
class of their minimum gap:

\medskip
\noindent\textbf{Category 1: Gap-2 transitions (minimum gap $\equiv 2 \pmod{3}$,
Class~II).}
The transitions $9\to1$, $7\to9$, $1\to3$ all have minimum gap~2. Gap~2 is the twin
prime gap: it carries the constitutional signature of Theorem~\ref{thm:signature}. The
Hardy--Littlewood singular series for gap~2 has a factor
\[
  \mathfrak{S}_2 = 2\prod_{p > 2}\frac{p(p-2)}{(p-1)^2} \approx 1.3203 > 1,
\]
which is the largest singular series value among even gaps. These transitions are
constitutionally enhanced: they represent pairs whose structure is closest to twin prime
pairs. The $9\to1$ transition is most enhanced (${\sim}41\%$) because its gap-2 landing
is arithmetically most accessible from that starting digit.

\medskip
\noindent\textbf{Category 2: Same-digit transitions (minimum gap $= 10$, $\equiv 1
\pmod{3}$, Class~I).}
Each self-transition ($1\to1$, $3\to3$, $7\to7$, $9\to9$) has minimum gap~10: a
prime must traverse an entire decimal cycle before returning to the same terminal digit.
The Hardy--Littlewood singular series for gap~10 satisfies $\mathfrak{S}_{10} < 1$.
These transitions are penalised by the geometric cost of the decimal cycle. This is the
direct cause of the diagonal suppression: it is not that primes ``avoid'' their own digit
by some mysterious repulsion, but that the minimum gap required to stay is the largest
of any transition, and the corresponding singular series is smallest.

\medskip
\noindent\textbf{Category 3: Gap-6 transitions (minimum gap $\equiv 0 \pmod{3}$,
Forbidden class).}
The transitions $1\to7$, $3\to9$, $7\to3$ have minimum gap~6. Gap~6 is in the Forbidden
ternary class: $6 \equiv 0 \pmod{3}$. The Hardy--Littlewood singular series for
Forbidden-class gaps has a factor
\[
  \chi_3(k) = \#\{j \in \{1,2\} : (j+k) \bmod 3 \in \{1,2\}\}
\]
which evaluates to zero when $k \equiv 0 \pmod{3}$, so $\mathfrak{S}_6 = 0$ for the
ternary factor at $p=3$. The probability mass that would have gone to gap-6 transitions
is redistributed to neighbouring gaps. This redistribution, not a constitutional
``preference,'' explains the apparent elevation of some transitions in this category.

\subsection{Formal Proposition}

\begin{proposition}[Min-Gap Ternary Transition Structure]
\label{prop:mingap}
Let $f(d_1, d_2)$ denote the observed transition probability from terminal digit $d_1$
to terminal digit $d_2$ among consecutive primes. Let $\mgap{d_1}{d_2}$ denote the
minimum possible prime gap for the transition. Then:
\begin{enumerate}[label=(\alph*)]
  \item $f(d_1, d_2)$ is \emph{enhanced} when $\mgap{d_1}{d_2} = 2$ (gap-2 transitions
    carry the twin prime constitutional signature);
  \item $f(d_1, d_2)$ is \emph{suppressed} when $d_1 = d_2$ (same-digit transitions
    require the maximum decimal gap of~10, penalised by the smallest singular series);
  \item $f(d_1, d_2)$ exhibits \emph{redistribution} when $\mgap{d_1}{d_2} \equiv 0
    \pmod{3}$ (the $p=3$ Euler factor of the Hardy--Littlewood singular series vanishes
    for Forbidden-class gaps, displacing probability mass to adjacent gaps).
\end{enumerate}
\end{proposition}

\begin{remark}
Proposition~\ref{prop:mingap} does not prove the Hardy--Littlewood conjecture---it
inherits the same unproven assumptions. What it provides is the mechanistic layer
\emph{beneath} that conjecture: the ternary class of the minimum gap is the reason the
singular series encodes the observed bias. The Hardy--Littlewood singular series is the
quantitative expression of ternary constitutional structure at the level of prime pairs.
\end{remark}

\begin{figure}[htbp]
\centering
\includegraphics[width=\textwidth]{fig2_mechanism_corrected}
\caption{Terminal digits $\{0$--$9\}$ coloured by ternary constitutional class (corrected
diagram). The four prime terminal digits $\{1, 3, 7, 9\}$ are shown bright; even and
composite terminal digits are faded. The red arcs mark the three gap-2 transitions
($9\to1$, $7\to9$, $1\to3$): these are constitutionally enhanced because minimum gap~2
carries the twin prime constitutional signature ($\equiv 2 \pmod{3}$, Class~II). The
dashed purple arc marks a gap-6 transition ($1\to7$), whose minimum gap is in the
Forbidden class ($6 \equiv 0 \pmod{3}$), causing redistribution rather than enhancement.
\emph{Note:} the structural driver is the ternary class of the minimum gap between
digits---not the ternary class of the terminal digit itself. No prime greater than~3 can
reside in the Forbidden class.}
\label{fig:mechanism}
\end{figure}

\subsection{Connection to the Hardy--Littlewood Singular Series}

The formal connection between ternary constitutional structure and the Hardy--Littlewood
singular series is through the $p = 3$ Euler factor. For an admissible $k$-tuple with
common difference $k$, define:
\[
  \chi_3(k) = \#\{j \in \{1, 2\} : (j + k) \bmod 3 \in \{1, 2\}\}.
\]
Then the singular series factor at $p = 3$ is:
\[
  \mathfrak{S}_3(k) = \frac{(1 - 1/3)^{-1}(1 - \chi_3(k)/\phi(3))}{1}
  = \begin{cases}
    \text{vanishes} & \text{if } k \equiv 0 \pmod{3}, \\
    \tfrac{1}{2} & \text{if } k \not\equiv 0 \pmod{3}.
  \end{cases}
\]
Constitutional Forcing provides the geometric first principles from which this singular
series structure emerges: the vanishing for $k \equiv 0 \pmod{3}$ is the algebraic fact
that gap-6 transitions map primes to the Forbidden class. The enhancement for gap-2
transitions is the constitutional signature of Theorem~\ref{thm:signature}. The
Hardy--Littlewood conjecture quantifies what constitutional structure constrains.

% ============================================================
\section{The Snail's Pace: Identifying the Convergence Rate}
\label{sec:snail}
% ============================================================

Lemke Oliver and Soundararajan noted that the bias diminishes as primes grow larger,
but ``at a snail's pace.'' This rate is not mysterious once the ternary framework is
in view.

The constitutional constraint operates at the level of residues modulo~3. As primes
grow, Theorem~\ref{thm:equilibrium} guarantees that $\pi_1(N)$ and $\pi_2(N)$ converge
to equal density. The deviation from the 50/50 split is governed by the
Siegel--Walfi{}sz theorem~\cite{Wal63}:
\begin{equation}
  \label{eq:sw}
  |\pi_1(N) - \pi_2(N)| = O\!\left(N \cdot \exp\!\left(-c\sqrt{\log N}\right)\right),
\end{equation}
for an absolute constant $c > 0$.

The exponential-of-square-root form in~\eqref{eq:sw} is famously slow
decay---much slower than any polynomial rate. This is precisely the snail's pace.

The digit-level bias tracks the constitutional-level equilibration. As the ternary
classes $\pi_1(N)/\pi(N)$ and $\pi_2(N)/\pi(N)$ converge to $1/2$, the surface-level
digit bias must diminish correspondingly, at exactly the rate specified by~\eqref{eq:sw}.
There is no mystery in the rate once one recognizes what is converging.

\begin{proposition}[Snail's Pace Identification]
\label{prop:snail}
The snail's pace of the Lemke Oliver--Soundararajan bias is the Siegel--Walfi{}sz
convergence rate of ternary class equilibration, projected onto decimal terminal digits.
The rate is not an empirical observation requiring separate explanation---it is the
canonical convergence rate of primes modulo small moduli, appearing here in a context
where it had not previously been recognised.
\end{proposition}

\begin{figure}[htbp]
\centering
\includegraphics[width=\textwidth]{fig3_snails_pace}
\caption{The Siegel--Walfi{}sz rate identifies the snail's pace. \emph{Left:} convergence
rate comparison on a log scale. The Siegel--Walfi{}sz envelope $e^{-c\sqrt{\log N}}$
decays far more slowly than any polynomial rate ($N^{-1/2}$, $N^{-1}$) or logarithmic
rate ($1/\log N$). The orange dashed line marks the scale of the Lemke
Oliver--Soundararajan study (first billion primes). \emph{Right:} the bias persists into
astronomical and trans-computational scales, confirming that the snail's pace
characterisation reflects a fundamental analytic fact, not a curiosity of the
computational range.}
\label{fig:snail}
\end{figure}

% ============================================================
\section{Validation, Context, and the Parity Barrier}
\label{sec:validation}
% ============================================================

\subsection{Base Independence}

The Lemke Oliver--Soundararajan bias holds across all number bases, not only base~10.
This is immediate from the ternary explanation: the mechanism operates in mod-3
arithmetic, which is independent of the presentation base. In any base~$b$, the
available prime terminal digits are those coprime to~$b$; each maps to a ternary class
modulo~3; and the minimum gaps between those digits determine the transition structure
through the same mechanism of Table~\ref{tab:mingaps}. The bias manifests differently
in each base because the minimum gap structure differs, but the underlying cause is
identical.

\begin{figure}[htbp]
\centering
\includegraphics[width=\textwidth]{fig4_base_independence}
\caption{Ternary constitutional class structure across number bases 6, 10, and~16. In
each base, the prime terminal digits (bright tiles) are those coprime to the base. The
distribution of those digits across the three ternary classes ($\equiv 0, 1, 2 \pmod{3}$)
varies by base, but the constitutional pressure mechanism is identical: the ternary class
of the minimum gap for each transition determines enhancement, suppression, or
redistribution. Base~6 has no Forbidden-class prime terminal digits; base~10 has two
(digits 3 and~9); base~16 has three (digits 3, 9, and~15). The bias projection differs;
the cause is identical.}
\label{fig:bases}
\end{figure}

\subsection{The Parity Barrier Context}

The ternary constitutional framework provides context for the longstanding parity barrier
in sieve theory. The Selberg parity barrier~\cite{Sel50} states that sieve methods
operating modulo~2 cannot distinguish between numbers with an odd versus even number
of prime factors. This structural ceiling stalled the bound on prime gaps at~246
(Maynard~\cite{May15}, Polymath8b).

The ternary approach is orthogonal to mod-2 arithmetic. Consider the ternary signatures
of consecutive prime gaps:

\begin{center}
\renewcommand{\arraystretch}{1.25}
\begin{tabular}{cccc}
\toprule
Gap size & $p \pmod{3}$ & $p + \text{gap} \pmod{3}$ & Constitutional? \\
\midrule
2  & 2 & 1 & \checkmark\ Unique Class~II $\to$ Class~I signature \\
4  & 2 & 0 (Forbidden) & $\times$\ Hits Forbidden class \\
6  & 2 & 2 (same class) & $\times$\ Same-class pairing \\
8  & 2 & 1 & $\sim$\ Similar to gap-2 but non-unique \\
\bottomrule
\end{tabular}
\end{center}

The gap-2 ternary signature $(2, 1)$ is the only even-gap signature that avoids the
Forbidden class for both members. Gap~4 immediately hits the Forbidden class. Gap~6
produces same-class pairing. This structural uniqueness is what allows the ternary method
to isolate twin primes specifically---the task that sieve methods operating mod~2 cannot
accomplish because modulo~2, all even gaps are identical.

\subsection{Relation to the Constitutional Forcing Programme}
\label{sec:relation}

This paper is a companion result to the Constitutional Forcing programme~\cite{Wakil26},
which establishes a level of distribution $\theta_W = 5/8$ for primes over cascade moduli
$q = 3^K$ via the Eisenstein integer geometry of $\mathbb{Z}[\omega]$.

The Lemke Oliver--Soundararajan explanation is independent of that result in the
following sense: it requires only the elementary constitutional constraint
(Definition~\ref{def:constraint}), the Siegel--Walfi{}sz theorem, and the minimum gap
class mapping. The full cascade moduli machinery is not needed to explain the digit bias.
The bias is a base-10 surface symptom of the same ternary structure that, pursued
deeper, yields the level-of-distribution result.

The relation is: the digit bias is what ternary constitutional pressure looks like at the
surface. The cascade moduli result is what it looks like at the analytic depth. Both are
projections of the same underlying structure.

% ============================================================
\section{Conclusion}
\label{sec:conclusion}
% ============================================================

The Lemke Oliver--Soundararajan digit bias was described when discovered as something
that ``every single person\ldots immediately went and wrote their own program to
check''~\cite{Gran16}. It has resisted mechanistic explanation for nine years.

We have shown that the bias is transparent once viewed in the correct coordinate system.
The primes do not ``know'' anything. They obey a structural constraint---the mod-3
constitutional requirement---that is elementary, provable in two lines, and that projects
onto base-10 terminal digits as the observed repulsion pattern. Three consequences:
\begin{enumerate}[label=(\roman*)]
  \item The diagonal suppression (primes avoiding their own terminal digit) is the
    structural penalty of same-digit transitions having the largest minimum gap (gap~10)
    and the smallest Hardy--Littlewood singular series value.
  \item The snail's pace of diminishment is the Siegel--Walfi{}sz rate of ternary class
    equilibration---one of the most fundamental quantities in analytic number theory,
    appearing here unrecognised.
  \item The base-independence of the bias is immediate: the mechanism is mod-3, not
    base-10.
\end{enumerate}

The corrected mechanism---min-gap ternary class structure rather than terminal-digit
ternary class---is stronger than the original claim. It correctly assigns the full $4 \times 4$
transition matrix to three categories (gap-2 enhanced, gap-6 redistributed, same-digit
suppressed) with no exceptions, and connects directly to the vanishing and non-vanishing
of the $p=3$ Euler factor in the Hardy--Littlewood singular series.

The deeper lesson is epistemological. The assumption of prime randomness has been so
foundational that the elementary mod-3 structure of primes was not examined as a source
of structural bias. The constitutional framework inverts this: accept the provable
structural constraints as the natural entry point, and the apparent mysteries dissolve.

When Freeman Dyson called the Montgomery--Dyson connection a miracle, he meant that
mathematics gave him no other word. The word for the Lemke Oliver--Soundararajan
bias, we submit, is: \emph{ternary}.

% ============================================================
% REFERENCES
% ============================================================
\bibliographystyle{plain}
\begin{thebibliography}{99}

\bibitem[LOS16]{LOS16}
R.~J. Lemke~Oliver and K.~Soundararajan,
\newblock Unexpected biases in the distribution of consecutive primes,
\newblock \emph{Proc. Natl. Acad. Sci. USA}, 113(31):E4446--E4454, 2016.
\newblock \href{https://arxiv.org/abs/1603.03720}{arXiv:1603.03720 [math.NT]}.

\bibitem[LOS17]{LOS17}
R.~J. Lemke~Oliver and K.~Soundararajan,
\newblock The distribution of consecutive prime biases and sums of sawtooth
  random variables,
\newblock \href{https://arxiv.org/abs/1709.06168}{arXiv:1709.06168 [math.NT]},
  2017.

\bibitem[HL23]{HL23}
G.~H. Hardy and J.~E. Littlewood,
\newblock Some problems of `Partitio numerorum'; III: on the expression of a
  number as a sum of primes,
\newblock \emph{Acta Math.}, 44:1--70, 1923.

\bibitem[Dir37]{Dir37}
P.~G.~L. Dirichlet,
\newblock Beweis des Satzes, dass jede unbegrenzte arithmetische Progression,
  deren erstes Glied und Differenz ganze Zahlen ohne gemeinschaftlichen Factor
  sind, unendlich viele Primzahlen enth\"alt,
\newblock \emph{Abhandlungen der K\"oniglich Preussischen Akademie der
  Wissenschaften}, pages 45--81, 1837.

\bibitem[Wal63]{Wal63}
A.~Walfi{}sz,
\newblock \emph{Weylsche Exponentialsummen in der neueren Zahlentheorie},
  Mathematische Forschungsberichte XV.
\newblock VEB Deutscher Verlag der Wissenschaften, Berlin, 1963.

\bibitem[Bom65]{Bom65}
E.~Bombieri,
\newblock On the large sieve,
\newblock \emph{Mathematika}, 12:201--225, 1965.

\bibitem[Zha14]{Zha14}
Y.~Zhang,
\newblock Bounded gaps between primes,
\newblock \emph{Ann. of Math.}, 179(3):1121--1174, 2014.

\bibitem[May15]{May15}
J.~Maynard,
\newblock Small gaps between primes,
\newblock \emph{Ann. of Math.}, 181(1):383--413, 2015.

\bibitem[Mon73]{Mon73}
H.~L. Montgomery,
\newblock The pair correlation of zeros of the zeta function,
\newblock In \emph{Analytic Number Theory}, Proc. Sympos. Pure Math.~24, pages
  181--193. Amer. Math. Soc., Providence, RI, 1973.

\bibitem[Sel50]{Sel50}
A.~Selberg,
\newblock An elementary proof of the prime-number theorem for arithmetic
  progressions,
\newblock \emph{Canad. J. Math.}, 2:66--78, 1950.

\bibitem[Pas25]{Pas25}
A.~Pascadi,
\newblock Primes with Beatty and Chebotarev conditions, and a note on Selberg's
  sieve,
\newblock \href{https://arxiv.org/abs/2505.00653}{arXiv:2505.00653v2
  [math.NT]}, 2025.

\bibitem[Wakil26]{Wakil26}
K.~Wakil,
\newblock Constitutional Forcing: A Wakil Programme,
\newblock ARC Institute of Knowware Technical Report, Calgary, Alberta, 2026.
\newblock Seven-paper series: Papers~00--06.

\bibitem[Wakil26b]{Wakil26b}
K.~Wakil,
\newblock Spherical Cow Philosophy: A Methodological Framework for Mathematical
  Discovery (Paper~02 of the Constitutional Forcing programme),
\newblock ARC Institute of Knowware, February 2026.

\bibitem[Gran16]{Gran16}
A.~Granville,
\newblock Quoted in ``Mathematicians Find a Peculiar Pattern in Primes,''
\newblock \emph{Science News}, March~15, 2016.

\bibitem[HY25]{HY25}
S.~A. Hartnoll and M.~Yang,
\newblock Prime numbers in quantum gravity near spacelike singularities,
\newblock Preprint, University of Cambridge, 2025.

\end{thebibliography}

\end{document}
